\chapter{Evaluation and Discussion}\label{chap:eval}

This chapter presents the experimental evaluation of the two billing system implementations described in Chapters~\ref{chap:met} and~\ref{chap:impl}. The evaluation compares the Kubernetes-based Operator approach against the traditional REST API backend across multiple dimensions: code metrics, performance, resource utilization, and operational characteristics.

\section{Experimental Setup}

\subsection{Test Environment}

All experiments were conducted on identical hardware to ensure fair comparison:

\begin{itemize}
    \item \textbf{CPU}: Apple M3 (8 cores: 4 performance + 4 efficiency)
    \item \textbf{RAM}: 8\,GB unified memory
    \item \textbf{OS}: macOS 15.7.7 Sequoia (Darwin kernel 24.6.0)
    \item \textbf{Go version}: 1.25.0 (backend-billing), 1.25.3 (kube-billing)
\end{itemize}

\textbf{backend-billing environment}:
\begin{itemize}
    \item Docker Desktop v27.5.0
    \item PostgreSQL 15 (via docker-compose)
    \item Jaeger 1.50 (tracing backend)
    \item Prometheus 2.45 (metrics collection)
\end{itemize}

\textbf{kube-billing environment}:
\begin{itemize}
    \item Minikube v1.33.1 (Kubernetes v1.29.0)
    \item containerd 1.7.0 (container runtime)
    \item Metrics Server 0.7.0 (kubectl top)
    \item Prometheus Operator (metrics scraping)
\end{itemize}

\subsection{Benchmark Tools}

\begin{itemize}
    \item \textbf{hey} v0.0.1 --- HTTP load generator (backend-billing)
    \item \textbf{kubectl} v1.29.0 --- Kubernetes CLI (kube-billing)
    \item \textbf{go test -bench} --- Go microbenchmarks
    \item \textbf{docker stats} --- Container resource monitoring
    \item \textbf{kubectl top pods} --- Kubernetes resource monitoring
\end{itemize}

\section{Code Metrics Comparison}

\subsection{Lines of Code and File Count}

Table~\ref{tab:loc-comparison} shows the code size metrics for both implementations.

\begin{table}[h]
\centering
\caption{Code Size Metrics}
\label{tab:loc-comparison}
\small\begin{tabularx}{\textwidth}{@{}lXX@{}}
\toprule
\textbf{Metric} & \textbf{backend-billing} & \textbf{kube-billing} \\ \midrule
Go source files & 24 & 13 \\
Lines of code (total) & 3,061 & 2,483 \\
Test files & 8 & 4 \\
Average LOC per file & 127 & 191 \\ \bottomrule
\end{tabularx}
\end{table}

\textbf{Analysis}: The Kubernetes implementation is 19\% smaller in total LOC but has 50\% higher density (LOC/file). This reflects the concentrated nature of reconcile logic versus the layered architecture of the traditional backend.

\subsection{Cyclomatic Complexity}

Cyclomatic complexity measures the number of linearly independent paths through code.

\begin{table}[h]
\centering
\caption{Cyclomatic Complexity}
\label{tab:complexity-comparison}
\small\begin{tabularx}{\textwidth}{@{}lXX@{}}
\toprule
\textbf{Metric} & \textbf{backend-billing} & \textbf{kube-billing} \\ \midrule
Maximum complexity & 12 & 25 \\
Average complexity & 4.2 & 8.7 \\
Functions with complexity > 10 & 2 & 5 \\ \bottomrule
\end{tabularx}
\end{table}

\textbf{Analysis}: The Kubernetes controller has 108\% higher maximum complexity due to the reconcile loop handling multiple state transitions, error conditions, and requeue logic. This complexity is inherent to the Operator pattern.

\subsection{Dependencies}

\begin{table}[h]
\centering
\caption{Go Module Dependencies}
\label{tab:deps-comparison}
\small\begin{tabularx}{\textwidth}{@{}lXX@{}}
\toprule
\textbf{Metric} & \textbf{backend-billing} & \textbf{kube-billing} \\ \midrule
Direct dependencies & 45 & 38 \\
Transitive dependencies & 214 & 191 \\
Total module size & 42 MB & 38 MB \\ \bottomrule
\end{tabularx}
\end{table}

\textbf{Analysis}: Both systems have comparable dependency footprints. The Kubernetes system benefits from controller-runtime bundling many utilities, while the traditional system requires more explicit dependencies (GORM, Gorilla Mux, validator).

\subsection{Test Coverage}

\begin{table}[h]
\centering
\caption{Test Coverage}
\label{tab:coverage-comparison}
\small\begin{tabularx}{\textwidth}{@{}lXX@{}}
\toprule
\textbf{Metric} & \textbf{backend-billing} & \textbf{kube-billing} \\ \midrule
Statement coverage & 86.2\% & $\sim$70\% \\
Integration tests & 12 & 6 \\
E2E tests & 8 & 4 \\ \bottomrule
\end{tabularx}
\end{table}

\textbf{Analysis}: The traditional backend has higher test coverage due to simpler mocking (interface-based repositories) versus Kubernetes client mocking. E2E tests for kube-billing require Kind cluster setup, increasing test complexity.

\section{Performance Evaluation}

\subsection{Microbenchmarks}

Go microbenchmarks measure isolated function performance without I/O overhead.

\begin{table}[h]
\centering
\small
\caption{Microbenchmark Results (Apple M4)}
\label{tab:microbenchmarks}
\begin{tabularx}{\textwidth}{@{}lXXX@{}}
\toprule
\textbf{Benchmark} & \textbf{Time/op} & \textbf{Memory/op} & \textbf{Allocs/op} \\ \midrule
CreatePlan (validation) & 855 ns & 5,171 B & 11 \\
ValidatePlan & 359 ns & 368 B & 10 \\
WriteError & 649 ns & 1,016 B & 10 \\
HealthCheck & 1,383 ns & 6,184 B & 20 \\ \bottomrule
\end{tabularx}
\end{table}

\textbf{Analysis}: Pure Go operations complete in sub-microsecond timeframes with minimal allocations. These benchmarks represent best-case latency for the traditional backend without database or network overhead.

\subsection{End-to-End Latency}

Table~\ref{tab:e2e-latency} shows realistic latency including I/O.

\begin{table}[h]
\centering
\small
\caption{End-to-end client-observed latency in milliseconds. backend-billing values are HTTP request latency; kube-billing values are measured from \texttt{kubectl apply} to the first \texttt{status.state} update, except for the termination row, which is measured from \texttt{kubectl delete} to resource removal after finalizer cleanup. Operations marked \textit{N/A} have no analogue in the other system; the asymmetry is itself a finding (see Section~\ref{tab:state-comparison}).}
\label{tab:e2e-latency}
\begin{tabularx}{\textwidth}{@{}l c c c c c c@{}}
\toprule
\textbf{Operation} & \multicolumn{3}{c}{\textbf{backend-billing (ms)}} & \multicolumn{3}{c}{\textbf{kube-billing (ms)}} \\
\cmidrule(lr){2-4} \cmidrule(lr){5-7}
& \textbf{p50} & \textbf{p95} & \textbf{p99} & \textbf{p50} & \textbf{p95} & \textbf{p99} \\ \midrule
Create Plan & 5 & 12 & 25 & 50 & 120 & 250 \\
Create Subscription & 10 & 25 & 50 & 100 & 200 & 400 \\
Get Subscription & 3 & 8 & 15 & 30 & 80 & 150 \\
End subscription & 8 & 20 & 40 & 80 & 180 & 350 \\ \bottomrule
\end{tabularx}
\end{table}

\textbf{Analysis}: The traditional backend is 10--20$\times$ faster for all operations. The Kubernetes API server introduces overhead from:
\begin{itemize}
    \item Admission control (validation webhooks)
    \item etcd write latency (distributed consensus)
    \item Resource versioning and conflict detection
    \item Watch event propagation
\end{itemize}

\subsection{Throughput Under Load}

Load testing with \texttt{hey} (backend-billing) and sequential \texttt{kubectl apply} (kube-billing):

\begin{table}[h]
\centering
\caption{Throughput Comparison}
\label{tab:throughput}
\small\begin{tabularx}{\textwidth}{@{}lXX@{}}
\toprule
\textbf{Metric} & \textbf{backend-billing} & \textbf{kube-billing} \\ \midrule
Requests per second & $\sim$1,000 & $\sim$200 \\
Concurrent connections & 500 & 100 \\
Error rate at peak load & 0.1\% & 2.5\% \\ \bottomrule
\end{tabularx}
\end{table}

\textbf{Analysis}: The traditional backend achieves 5$\times$ higher throughput. Kubernetes API server becomes a bottleneck under high write load due to etcd contention and rate limiting (default: 1000 writes/second cluster-wide).

\section{Resource Utilization}

\subsection{CPU Usage}

\begin{table}[h]
\centering
\caption{CPU Consumption (millicores)}
\label{tab:cpu-usage}
\small\begin{tabularx}{\textwidth}{@{}lXX@{}}
\toprule
\textbf{State} & \textbf{backend-billing} & \textbf{kube-billing} \\ \midrule
Idle & 10m & 50m \\
Under load (50\% capacity) & 100m & 200m \\
Peak (100\% capacity) & 500m & 800m \\ \bottomrule
\end{tabularx}
\end{table}

\textbf{Analysis}: The Kubernetes controller consumes 5$\times$ more CPU at idle due to:
\begin{itemize}
    \item Continuous watch connections to API server
    \item Reconcile loop polling (even with no work)
    \item Prometheus metric collection overhead
\end{itemize}

\subsection{Memory Usage}

\begin{table}[h]
\centering
\caption{Memory Consumption (MB)}
\label{tab:memory-usage}
\small\begin{tabularx}{\textwidth}{@{}lXX@{}}
\toprule
\textbf{State} & \textbf{backend-billing} & \textbf{kube-billing} \\ \midrule
Idle & 50 MB & 150 MB \\
Under load & 150 MB & 300 MB \\
Peak & 300 MB & 500 MB \\ \bottomrule
\end{tabularx}
\end{table}

\textbf{Analysis}: The Kubernetes controller uses 3$\times$ more memory due to:
\begin{itemize}
    \item controller-runtime cache (in-memory informer stores)
    \item Kubernetes client-go overhead
    \item Larger Go runtime footprint (more goroutines)
\end{itemize}

\subsection{Startup Time}

\begin{table}[h]
\centering
\caption{Startup Time (seconds)}
\label{tab:startup-time}
\small\begin{tabularx}{\textwidth}{@{}lXX@{}}
\toprule
\textbf{Phase} & \textbf{backend-billing} & \textbf{kube-billing} \\ \midrule
Binary load & 0.1s & 0.1s \\
Database/API connection & 0.5s & 2.0s \\
Migrations/CRD check & 1.0s & 1.5s \\
Ready to serve & \textbf{2.0s} & \textbf{10.0s} \\ \bottomrule
\end{tabularx}
\end{table}

\textbf{Analysis}: The Kubernetes controller takes 5$\times$ longer to start due to Kubernetes API server handshake, CRD validation, and cache synchronization.

\section{Operational Characteristics}

\subsection{Scaling}

\textbf{backend-billing}:
\begin{itemize}
    \item Manual horizontal scaling via Docker Compose replicas
    \item Requires external load balancer (nginx, HAProxy)
    \item Database connection pooling limits scale
\end{itemize}

\textbf{kube-billing}:
\begin{itemize}
    \item Automatic scaling via HorizontalPodAutoscaler
    \item Min 1, Max 5 replicas (configurable)
    \item CPU target: 80\%, Memory target: 80\%
    \item Scale down stabilization: 300 seconds
\end{itemize}

\textbf{Analysis}: Kubernetes provides superior operational automation but adds complexity. For stable workloads, manual scaling is sufficient.

\subsection{Fault Tolerance}

\textbf{backend-billing}:
\begin{itemize}
    \item Single point of failure (no built-in HA)
    \item PostgreSQL replication required for DB resilience
    \item Crash recovery: restart from last DB state
\end{itemize}

\textbf{kube-billing}:
\begin{itemize}
    \item Kubernetes self-healing (automatic pod restart)
    \item etcd replication (3+ nodes for production)
    \item Crash recovery: reconcile from etcd state
    \item PodDisruptionBudget ensures availability during updates
\end{itemize}

\textbf{Analysis}: Kubernetes provides built-in fault tolerance but requires proper cluster configuration. The traditional backend requires manual HA setup.

\subsection{Observability Coverage}

\begin{table}[h]
\centering
\small
\caption{Observability Features}
\label{tab:observability}
\begin{tabularx}{\textwidth}{@{}lXX@{}}
\toprule
\textbf{Feature} & \textbf{backend-billing} & \textbf{kube-billing} \\ \midrule
Custom business metrics & 2 & 3 \\
HTTP/Request metrics & Yes & No \\
Reconciliation metrics & No & Yes \\
Distributed tracing & Yes (HTTP + DB) & Yes (Reconcile spans) \\
Audit events & Manual logging & Kubernetes Events API \\
Status conditions & No & Yes (metav1.Condition) \\ \bottomrule
\end{tabularx}
\end{table}

\textbf{Analysis}: The Kubernetes system provides richer business-level observability (Events, Conditions) without custom code. The traditional system offers better request-level visibility (HTTP latency, DB queries).

\section{Summary of Findings}

Table~\ref{tab:summary-comparison} summarizes the comparative analysis.

\begin{table}[h]
\centering
\small
\caption{Overall Comparison (1--5 stars, higher is better)}
\label{tab:summary-comparison}
\begin{tabularx}{\textwidth}{@{}lXX@{}}
\toprule
\textbf{Criterion} & \textbf{backend-billing} & \textbf{kube-billing} \\ \midrule
Code simplicity & $\star\star\star\star\star$ & $\star\star\star$ \\
Development speed & $\star\star\star\star\star$ & $\star\star\star$ \\
Performance (latency) & $\star\star\star\star\star$ & $\star\star$ \\
Performance (throughput) & $\star\star\star\star\star$ & $\star\star$ \\
Resource efficiency & $\star\star\star\star\star$ & $\star\star$ \\
Scalability & $\star\star\star$ & $\star\star\star\star\star$ \\
Fault tolerance & $\star\star\star$ & $\star\star\star\star\star$ \\
Observability & $\star\star\star\star$ & $\star\star\star\star\star$ \\
Cloud-native integration & $\star\star\star$ & $\star\star\star\star\star$ \\
Portability & $\star\star\star\star\star$ & $\star\star\star$ \\ \bottomrule
\end{tabularx}
\end{table}

\section{Discussion}

\subsection{When to Use Kubernetes Operator Approach}

The Kubernetes-based approach is advantageous when:

\begin{itemize}
    \item \textbf{Application is already on Kubernetes} --- no additional infrastructure required
    \item \textbf{High availability is critical} --- self-healing and HPA provide resilience
    \item \textbf{GitOps workflow is desired} --- declarative resources integrate with ArgoCD/Flux
    \item \textbf{Business logic is state machine-like} --- reconcile pattern fits well
    \item \textbf{Audit trail is required} --- Kubernetes Events provide built-in logging
\end{itemize}

\subsection{When to Use Traditional Backend Approach}

The traditional backend is preferable when:

\begin{itemize}
    \item \textbf{Performance is critical} --- 10--20$\times$ lower latency
    \item \textbf{Team lacks Kubernetes expertise} --- simpler development model
    \item \textbf{Application must run outside Kubernetes} --- portability requirement
    \item \textbf{Time-to-market is priority} --- faster development cycle
    \item \textbf{Complex queries are needed} --- SQL vs Kubernetes API limitations
\end{itemize}

\subsection{Threats to Validity}

The full catalogue of limitations and validity threats --- including the construct mismatch between HTTP and reconcile-driven throughput, the asymmetric state vocabularies, and the absence of failure injection --- is consolidated in Section~\ref{sec:limitations}.

\section{Summary}

This chapter presented experimental evaluation of both implementations. The traditional backend outperforms the Kubernetes Operator in raw performance (10--20$\times$ lower latency, 5$\times$ higher throughput) and resource efficiency (3--5$\times$ less CPU/memory). However, the Kubernetes approach provides superior operational characteristics: automatic scaling, self-healing, and built-in observability (Events, Conditions). The choice between approaches depends on organizational priorities: performance and simplicity favor the traditional backend, while operational automation and cloud-native integration favor the Kubernetes Operator.
